% GENERATED FILE. Edit canonical lessons, then run npm run generate:books.
% canonical-source-hash: fnv1a64:d1696fce0719e817
% canonical-lessons: MW-C210-parakhno, MW-C210-bhat, MW-C210-namuno, MW-C210-nakli, MW-C210-chunno

\chapter{Choosing the Right One}
\label{ch:mw-choosing-the-right-one}
\hlchaptermodality{210}

\begin{chapteropening}
\textbf{By the end of this chapter:} \emph{I can check what I buy and pick one out.}

Complete the last lesson of chapter 210: I can check what I buy and pick one out.
\end{chapteropening}

\section[\texorpdfstring{\mw{परखणो} (parakhṇo)}{परखणो (parakhṇo)}]{\mw{परखणो} (parakhṇo) — to test; to check}

\label{lesson:MW-C210-parakhno}



\begin{quote}
Before the new one: say the Marwadi for a deal, then the Marwadi for to order.
\end{quote}

\subsection*{You'll want to know: \mw{परखणो}}
\textbf{\mw{परखणो}} — \emph{parakhṇo} — \textquotedblleft{}to test; to check\textquotedblright{}.

\begin{quote}
\mw{मां} \mw{कपड़ा} \mw{परखै} \mw{है।}
\end{quote}

\begin{quote}
\emph{mā̃ kapṛā parkhai hai.}
\end{quote}

\begin{quote}
Mother checks the cloth.
\end{quote}

Past (I …): \mw{म्हैं} \mw{कपड़ा} \mw{परख्या}, \emph{mhaĩ kapṛā parakhyā}.

\begin{cousinweb}
\textbf{\mw{परखणो}} (\emph{parakhṇo}) is inherited from Sanskrit \emph{parīkṣate}, \textquotedblleft{}observes, examines\textquotedblright{}; CDIAL lists an Old Marwari form beside Hindi \emph{parakhnā}, \textquotedblleft{}to test\textquotedblright{} (CDIAL 7904).
\end{cousinweb}

\begin{grammarlens}[title={What you've built}]
One more word for this chapter.
\end{grammarlens}

\subsection*{Your turn}
\begin{itemize}
\raggedright
  \item \emph{Say it:} \emph{parakhṇo}
  \item \emph{Say it:} \emph{parakhṇo}, once more
  \item \emph{Say it:} \emph{mā̃ kapṛā parkhai hai.} — Mother checks the cloth.
  \item \emph{Recall:} say the Marwadi for a deal, then the Marwadi for to order, then say \emph{parakhṇo} again
\end{itemize}

\begin{tcolorbox}[breakable,colback=teal!4,colframe=teal!35!black,title={Before you move on}]
What does \mw{परखणो} mean? (\textquotedblleft{}To test; to check\textquotedblright{}.) Say it once more.
\end{tcolorbox}
\section[\texorpdfstring{\mw{भांत} (bhā̃t)}{भांत (bhā̃t)}]{\mw{भांत} (bhā̃t) — a sort; a kind}

\label{lesson:MW-C210-bhat}



\begin{quote}
Before the new one: say the Marwadi for to order, then the Marwadi for to test.
\end{quote}

\subsection*{You'll want to know: \mw{भांत}}
\textbf{\mw{भांत}} — \emph{bhā̃t} — \textquotedblleft{}a sort; a kind\textquotedblright{}.

\begin{quote}
\mw{दूजी} \mw{भांत} \mw{दिखावो।}
\end{quote}

\begin{quote}
\emph{dūjī bhā̃t dikhāvo.}
\end{quote}

\begin{quote}
Show me another kind.
\end{quote}

A feminine noun.

\begin{cousinweb}
\textbf{\mw{भांत}} (\emph{bhā̃t}) is inherited from Sanskrit \emph{bhakti}, \textquotedblleft{}a share; an arrangement, a manner\textquotedblright{}; CDIAL lists Old Marwari \emph{bhā̃ti}, \textquotedblleft{}manner, kind\textquotedblright{} (CDIAL 9338).
\end{cousinweb}

\begin{grammarlens}[title={What you've built}]
One more word for this chapter.
\end{grammarlens}

\subsection*{Your turn}
\begin{itemize}
\raggedright
  \item \emph{Say it:} \emph{bhā̃t}
  \item \emph{Say it:} \emph{bhā̃t}, once more
  \item \emph{Say it:} \emph{dūjī bhā̃t dikhāvo.} — Show me another kind.
  \item \emph{Recall:} say the Marwadi for to order, then the Marwadi for to test, then say \emph{bhā̃t} again
\end{itemize}

\begin{tcolorbox}[breakable,colback=teal!4,colframe=teal!35!black,title={Before you move on}]
What does \mw{भांत} mean? (\textquotedblleft{}A sort; a kind\textquotedblright{}.) Say it once more.
\end{tcolorbox}
\section[\texorpdfstring{\mw{नमूनो} (namūno)}{नमूनो (namūno)}]{\mw{नमूनो} (namūno) — a sample}

\label{lesson:MW-C210-namuno}



\begin{quote}
Before the new one: say the Marwadi for to test, then the Marwadi for a sort.
\end{quote}

\subsection*{You'll want to know: \mw{नमूनो}}
\textbf{\mw{नमूनो}} — \emph{namūno} — \textquotedblleft{}a sample\textquotedblright{}.

\begin{quote}
\mw{ये} \mw{नमूनो} \mw{आछो} \mw{है।}
\end{quote}

\begin{quote}
\emph{ye namūno āchho hai.}
\end{quote}

\begin{quote}
This sample is good.
\end{quote}

A masculine noun.

\begin{cousinweb}
\textbf{\mw{नमूनो}} (\emph{namūno}) is usually traced to Persian \emph{namūna}, \textquotedblleft{}a sample, a model\textquotedblright{}.
\end{cousinweb}

\begin{grammarlens}[title={What you've built}]
One more word for this chapter.
\end{grammarlens}

\subsection*{Your turn}
\begin{itemize}
\raggedright
  \item \emph{Say it:} \emph{namūno}
  \item \emph{Say it:} \emph{namūno}, once more
  \item \emph{Say it:} \emph{ye namūno āchho hai.} — This sample is good.
  \item \emph{Recall:} say the Marwadi for to test, then the Marwadi for a sort, then say \emph{namūno} again
\end{itemize}

\begin{tcolorbox}[breakable,colback=teal!4,colframe=teal!35!black,title={Before you move on}]
What does \mw{नमूनो} mean? (\textquotedblleft{}A sample\textquotedblright{}.) Say it once more.
\end{tcolorbox}
\section[\texorpdfstring{\mw{नकली} (naklī)}{नकली (naklī)}]{\mw{नकली} (naklī) — fake}

\label{lesson:MW-C210-nakli}



\begin{quote}
Before the new one: say the Marwadi for a sort, then the Marwadi for a sample.
\end{quote}

\subsection*{You'll want to know: \mw{नकली}}
\textbf{\mw{नकली}} — \emph{naklī} — \textquotedblleft{}fake\textquotedblright{}.

\begin{quote}
\mw{ये} \mw{घड़ी} \mw{नकली} \mw{है।}
\end{quote}

\begin{quote}
\emph{ye ghaṛī naklī hai.}
\end{quote}

\begin{quote}
This watch is a fake.
\end{quote}

\begin{cousinweb}
\textbf{\mw{नकली}} (\emph{naklī}) is usually traced, through Persian, to Arabic \emph{naql}, \textquotedblleft{}a copy\textquotedblright{}.
\end{cousinweb}

\begin{grammarlens}[title={What you've built}]
One more word for this chapter.
\end{grammarlens}

\subsection*{Your turn}
\begin{itemize}
\raggedright
  \item \emph{Say it:} \emph{naklī}
  \item \emph{Say it:} \emph{naklī}, once more
  \item \emph{Say it:} \emph{ye ghaṛī naklī hai.} — This watch is a fake.
  \item \emph{Recall:} say the Marwadi for a sort, then the Marwadi for a sample, then say \emph{naklī} again
\end{itemize}

\begin{tcolorbox}[breakable,colback=teal!4,colframe=teal!35!black,title={Before you move on}]
What does \mw{नकली} mean? (\textquotedblleft{}Fake\textquotedblright{}.) Say it once more.
\end{tcolorbox}
\section[\texorpdfstring{\mw{चुणणो} (chuṇṇo)}{चुणणो (chuṇṇo)}]{\mw{चुणणो} (chuṇṇo) — to pick out; to choose}

\label{lesson:MW-C210-chunno}



\begin{quote}
Before the new one: say the Marwadi for a sample, then the Marwadi for fake.
\end{quote}

\subsection*{You'll want to know: \mw{चुणणो}}
\textbf{\mw{चुणणो}} — \emph{chuṇṇo} — \textquotedblleft{}to pick out; to choose\textquotedblright{}.

\begin{quote}
\mw{हूं} \mw{चोखी} \mw{जूती} \mw{चुणूं} \mw{हूं।}
\end{quote}

\begin{quote}
\emph{hū̃ chokhī jūtī chuṇū̃ hū̃.}
\end{quote}

\begin{quote}
I pick out a good pair of shoes.
\end{quote}

Past (I …): \mw{म्हैं} \mw{लाल} \mw{टोपी} \mw{चुणी}, \emph{mhaĩ lāl ṭopī chuṇī}.

\begin{cousinweb}
\textbf{\mw{चुणणो}} (\emph{chuṇṇo}) is inherited from Sanskrit \emph{cinoti}, \textquotedblleft{}heaps up, collects\textquotedblright{}. CDIAL gives no Marwari form in this entry, but lists Hindi \emph{cunnā}, \textquotedblleft{}to pick up, gather\textquotedblright{} (CDIAL 4814); picking out the best is probably a later step from gathering.
\end{cousinweb}

\begin{grammarlens}[title={What you've built}]
One more word for this chapter.
\end{grammarlens}

\subsection*{Your turn}
\begin{itemize}
\raggedright
  \item \emph{Say it:} \emph{chuṇṇo}
  \item \emph{Say it:} \emph{chuṇṇo}, once more
  \item \emph{Say it:} \emph{hū̃ chokhī jūtī chuṇū̃ hū̃.} — I pick out a good pair of shoes.
  \item \emph{Recall:} say the Marwadi for a sample, then the Marwadi for fake, then say \emph{chuṇṇo} again
\end{itemize}

\begin{tcolorbox}[breakable,colback=teal!4,colframe=teal!35!black,title={Before you move on}]
What does \mw{चुणणो} mean? (\textquotedblleft{}To pick out; to choose\textquotedblright{}.) Say it once more.
\end{tcolorbox}
